% at_commands.tex - the AT command set of D-Modem's modem core
% (core/at.cpp), as slmodemd, libdmodem and the PBX modules have it.
% Built with docs/at/build.sh into docs/at_commands.pdf.
\documentclass[a4paper,10pt]{article}

\usepackage[T1]{fontenc}
\usepackage[utf8]{inputenc}
\usepackage{lmodern}
\usepackage{microtype}
\usepackage[margin=2.2cm]{geometry}
\usepackage{booktabs}
\usepackage{longtable}
\usepackage{array}
\usepackage{enumitem}
\usepackage{xcolor}
\usepackage[hidelinks]{hyperref}

\setlist{nosep,leftmargin=*}
\setlength{\parindent}{0pt}
\setlength{\parskip}{0.6ex}

% a command or what the modem says, as typed
\newcommand{\cmd}[1]{\texttt{#1}}
% the backslash, as in AT\N
\newcommand{\bs}{\textbackslash}
% a parameter
\newcommand{\p}[1]{\textmd{\textit{#1}}}

% one command: its syntax, then what it does
\newenvironment{command}[1]
  {\par\medskip\noindent\textbf{\cmd{#1}}\par\nopagebreak
   \begin{list}{}{\leftmargin=1.5em\topsep=0pt}\item[]}
  {\end{list}}

\title{D-Modem AT Command Reference}
\author{The modem core of slmodemd, libdmodem and the PBX modules}
\date{October 2026}

\begin{document}
\maketitle
\tableofcontents

\section{Introduction}

D-Modem's modem core (\cmd{core/} in the source) takes the AT commands of
every D-Modem program:
\begin{itemize}
\item \textbf{slmodemd}, the modem on a pseudo terminal
  (\cmd{/dev/ttySL0} links to it), its line a SIP call;
\item \textbf{libdmodem} (\cmd{dmodem.h}), the modem as a library: the
  program that embeds it sends the commands;
\item the \textbf{Asterisk and FreeSWITCH modules} (\cmd{app\_dmodem},
  \cmd{mod\_dmodem}), whose profiles give the commands each call starts
  with.
\end{itemize}

The commands, and the way a command line is read, are those of Smart
Link's slmodemd (Debian's \cmd{sl-modem-daemon}), whose ST56 modems'
reference is kept as \cmd{docs/legacy/ST56ATCommands.pdf}.
Section~\ref{sec:differences} lists where D-Modem differs from it, and
section~\ref{sec:sip} what a SIP line changes in slmodemd.

This reference follows the core as it is: every command, parameter range,
default and output here is the core's own, the outputs as slmodemd
prints them.  A setting that is kept but changes nothing says so.

\paragraph{Notation.} \cmd{AT} commands are shown in capitals; the modem
takes either case.  \p{n} is a decimal number, \p{h} a hexadecimal one;
\verb|[ ]| marks what may be left out.  A setting's default is the one
\cmd{ATZ} and the start of the program give.

\section{Command lines}

\begin{itemize}
\item A command line begins with \cmd{AT} and ends with the character of
  S3 (carriage return) or S4 (line feed).  Blanks and tabs in it are
  dropped before it is read.
\item The character of S5 (backspace) takes back the last one typed.
\item A line holds 63 characters at most; a longer one is ERROR.
\item With echo on (\cmd{E1}), what the DTE types comes back as it is
  typed.
\item A line that does not begin with \cmd{AT} is not a command: with
  echo on it is blanked out on the DTE's screen.
\item \cmd{A/} (no line end needed) runs the last command line again.
\end{itemize}

The commands of a line run in turn.  The first that fails ends the line
with ERROR; otherwise OK follows the last one.  \cmd{A}, \cmd{D} and
\cmd{O} end the line and give their result later (CONNECT, NO CARRIER and
so on), and \cmd{Z} ends it with OK: whatever follows them on the line is
not read.  A command that reports something (\cmd{I}, \cmd{S}\p{n}\cmd{?},
\cmd{+MS?} \ldots) prints it on lines of its own before the result.

A character from the DTE while a call is being set up (dialing,
answering, the modems' training) hangs up.

\section{Result codes}

\begin{center}
\begin{tabular}{cll}
\toprule
\cmd{V0} & \cmd{V1} & \\
\midrule
0 & OK & the command line ran \\
1 & CONNECT [\p{rate}] & the call is up, or a fax command's data begins \\
2 & RING & a call rings (slmodemd: every 3 s, section~\ref{sec:sip}) \\
3 & NO CARRIER & no call came up, or the call is over \\
4 & ERROR & a command failed, or a line too long \\
6 & NO DIALTONE & no dial tone (\cmd{X2}, \cmd{X4}) \\
7 & BUSY & a busy tone \\
8 & NO ANSWER & no answer within S7 seconds of dialing \\
\bottomrule
\end{tabular}
\end{center}

\cmd{Q1} turns them off, \cmd{V0} gives the numbers, each on a line of
its own.  With \cmd{V1} a data call's CONNECT gives the rate the modem
receives at, in bit/s (\cmd{CONNECT 33600}); with \cmd{X0}, in fax and
in voice mode it is CONNECT alone.

Before CONNECT, \cmd{\&A} (S70) adds lines on the connection (verbose
results only), each bit one of them:

\begin{center}
\begin{tabular}{cll}
\toprule
bit & line & \\
\midrule
0 (1) & \cmd{Modulation: }\p{name} & the modulation's name (\cmd{V92},
  \cmd{V34}, \cmd{V32bis}, \ldots) \\
1 (2) & \cmd{Protocol: LAPM} or \cmd{NONE} & V.42 error correction \\
0 and 1 (3) & \cmd{Compression: V44}, \cmd{V42bis} or \cmd{NONE} &
  what V.42 agreed \\
2 (4) & \cmd{TxRate: }\p{rate} & the rate the modem sends at \\
\bottomrule
\end{tabular}
\end{center}

\section{Basic commands}

\begin{command}{ATA}
Answers: off hook, and the answering modem's tones after S92 seconds.
In slmodemd the call that rings (\cmd{RING}) is answered.  Its result is
CONNECT, or NO CARRIER.
\end{command}

\begin{command}{ATB}
V.22 at 1200 bit/s only, without automode: \cmd{AT+MS=22,0,1200,1200}.
(Not V.250's choice between the ITU and Bell modulations: those are
\cmd{+MS}'s.)
\end{command}

\begin{command}{ATD\p{string}}
Dials \p{string} and calls.  The dialer sends the digits and
\cmd{* \# A B C D} as DTMF (S11 ms each, as long between them); a comma
pauses S8 seconds, \cmd{W} two; other characters are skipped.  \cmd{T}
and \cmd{P} choose tone or pulse dialing (S16) for the rest of the call.
\cmd{ATDL} dials the last number again (\cmd{Dialing }\p{number}\cmd{...}
first).  With \cmd{X2} or \cmd{X4} the modem waits for a dial tone first.
The result is CONNECT, BUSY, NO ANSWER, NO DIALTONE or NO CARRIER.  In
slmodemd the dial string is the SIP call's number
(section~\ref{sec:sip}).
\end{command}

\begin{command}{ATE[\p{n}]}
Echo of the DTE's characters in command mode: \cmd{E0} off, \cmd{E1} on
(S13; default 1).
\end{command}

\begin{command}{ATH[0], ATH1}
\cmd{H} or \cmd{H0} hangs up (on hook); \cmd{H1} goes off hook.
\end{command}

\begin{command}{ATI[\p{n}]}
Information:
\begin{center}
\begin{tabular}{cl}
\toprule
\cmd{I0} & the product: \cmd{SmartLink Soft Modem} \\
\cmd{I1}, \cmd{I2} & the product and version (\cmd{SmartLink Soft Modem,
  2.9.11}), then \cmd{Smart Link Ltd.} \\
\cmd{I3} & the device's name, then the driver's (slmodemd:
  \cmd{/dev/slamr0}, \cmd{socket driver}) \\
\cmd{I4} & S0 to S23, eight a line (\cmd{s00=000 s01=000 \ldots}) \\
\cmd{I5}, \cmd{I6} & \cmd{Stored Profile 0:}, \cmd{1:} (there are none) \\
\cmd{I7} & the country: \cmd{Country: USA} \\
\bottomrule
\end{tabular}
\end{center}
Any other number is ERROR.
\end{command}

\begin{command}{ATL[\p{n}], ATM[\p{n}]}
The speaker's volume, \cmd{L0} to \cmd{L3} (S30; default 3), and when it
is on, \cmd{M0} never, \cmd{M1} until the call is up, \cmd{M2} always
(S29; default 1).  A SIP line has no speaker: kept for the commands
that set them.
\end{command}

\begin{command}{ATN[\p{n}]}
Automode (S31; default 1): \cmd{N1} lets the modems find their best
common modulation from the one \cmd{+MS} chose down, \cmd{N0} keeps to
it.  The same as \cmd{+MS}'s second parameter.
\end{command}

\begin{command}{ATO}
Back online after the escape (section~\ref{sec:online}): CONNECT.  ERROR
when there is no call to go back to.
\end{command}

\begin{command}{ATP, ATT}
Pulse or tone dialing (S16; default tone).
\end{command}

\begin{command}{ATQ[\p{n}]}
\cmd{Q0} gives the result codes, \cmd{Q1} keeps them back (S14; default
0).
\end{command}

\begin{command}{ATS\p{n}?, ATS\p{n}=\p{v}}
Reads S-register \p{n} (0 to 255), in decimal on a line of its own, or
sets it to \p{v} (8 bits: \p{v} modulo 256).  A number over 255 is
ERROR.  Section~\ref{sec:sregs} lists them.
\end{command}

\begin{command}{ATV[\p{n}]}
\cmd{V0} gives the result codes as numbers, \cmd{V1} as words (S15;
default 1).
\end{command}

\begin{command}{ATX[\p{n}]}
Results and dial tone, \cmd{X0} to \cmd{X4} (S56; default 4):
\begin{center}
\begin{tabular}{cll}
\toprule
& CONNECT & before dialing \\
\midrule
\cmd{X0} & without the rate & no wait for a dial tone \\
\cmd{X1}, \cmd{X3} & with the rate & no wait for a dial tone \\
\cmd{X2}, \cmd{X4} & with the rate & a dial tone of 1 s, NO DIALTONE if
  none in S6 seconds (2 s at least) \\
\bottomrule
\end{tabular}
\end{center}
BUSY (four cadences of 0.2 to 0.8 s) and NO ANSWER come in all of them.
A SIP call has no dial tone: slmodemd wants \cmd{ATX3}.
\end{command}

\begin{command}{ATY[\p{n}]}
\cmd{Y0} to \cmd{Y3}: kept in S161, nothing more.
\end{command}

\begin{command}{ATZ[\p{n}]}
\cmd{Z} or \cmd{Z0} to \cmd{Z3}: hangs up, and the S-registers,
\cmd{+MS}, modem-on-hold and the data mode (\cmd{+FCLASS=0}) as at the
start; the country stays.  There are no stored profiles: the number
changes nothing.  OK, and the rest of the line is not read.
\end{command}

\section{Ampersand, backslash and percent commands}

\begin{command}{AT\&A[\p{n}]}
\cmd{\&A0} to \cmd{\&A7}: the lines before CONNECT (S70; default 0), see
the result codes.
\end{command}

\begin{command}{AT\&C[\p{n}], AT\&E[\p{n}], AT\&H[\p{n}]}
Data carrier detect (\cmd{\&C0}, \cmd{\&C1}; S60), the rate CONNECT
gives (\cmd{\&E0}, \cmd{\&E1}; S71) and flow control (\cmd{\&H0},
\cmd{\&H1}; S62): kept, but they change nothing (a pty has no carrier
line, CONNECT always gives the receive rate, and flow control is the
pty's).  Default 0.
\end{command}

\begin{command}{AT\&D[\p{n}], AT\&F[\p{n}], AT\&R[\p{n}], AT\&S[\p{n}]}
Taken, with or without a digit, and nothing done (\cmd{ATZ} gives the
defaults).
\end{command}

\begin{command}{AT\&K[\p{n}], AT\%C[\p{n}]}
Data compression (S104; default 3): \cmd{0} none, \cmd{1} to \cmd{3}
V.44 or V.42bis both ways, as the other modem agrees (V.42 offers both).
It needs V.42 (\cmd{\bs N}).
\end{command}

\begin{command}{AT\&P[\p{n}]}
\cmd{\&P0} to \cmd{\&P3}: the pulse dialing ratio, kept in S28.
\end{command}

\begin{command}{AT\&V, AT\&W}
ERROR: there are no stored profiles to show or write.
\end{command}

\begin{command}{AT\bs A[\p{n}]}
\cmd{\bs A0} to \cmd{\bs A3}, MNP's block size: taken (there is no
MNP); \cmd{\bs A4} and up are ERROR.
\end{command}

\begin{command}{AT\bs N[\p{n}]}
Error correction (S103; default 1): \cmd{\bs N0} none; \cmd{\bs N1} to
\cmd{\bs N5} V.42 (LAPM) when the other modem has it, a call without
when it does not (V.42's detection phase decides).  There is no MNP:
an MNP modem that has V.42 connects with LAPM, one with MNP only
without error correction or not at all.
\end{command}

\begin{command}{AT\%E[\p{n}]}
\cmd{\%E0} to \cmd{\%E2}, line quality monitoring: kept in S39 (the
modulations retrain and renegotiate their rate by themselves).
\end{command}

\begin{command}{AT*B}
The numbers barred from dialing: none are kept; OK.
\end{command}

\begin{command}{AT\#CID=\p{n}, AT\#CID?, AT\#CID=?}
Caller ID between the first two rings: \cmd{0} off, \cmd{1} or \cmd{2}
on (\cmd{=?} gives \cmd{0,1,2}).  A SIP call rings without it.
\end{command}

\section{Extended commands}

\subsection{Modulation: +MS}

\begin{command}{AT+MS=[\p{mod}][,[\p{automode}][,[\p{min}][,[\p{max}]]]]}
The modulation the modem calls or answers with, whether it may fall
back (\p{automode} 1, as \cmd{N1}; S31) and the rates it may connect at,
\p{min} to \p{max} bit/s (300 to 56000, \p{min} no more than \p{max}).
A parameter left out keeps its value: \cmd{AT+MS=,0} only turns
automode off.  The default is \cmd{92,1,300,56000}.

\begin{center}
\begin{tabular}{rll}
\toprule
\p{mod} & & \\
\midrule
21 & ITU V.21 & 300 bit/s FSK \\
23 & ITU V.23 & 1200/75 bit/s FSK \\
103 & Bell 103 & 300 bit/s FSK \\
22 & ITU V.22 & 1200 bit/s DPSK \\
212 & Bell 212A & 1200 bit/s DPSK \\
122 & ITU V.22bis & 2400 bit/s, V.22 as its fallback \\
32 & ITU V.32 & 4800 and 9600 bit/s \\
132 & ITU V.32bis & 4800 to 14400 bit/s \\
34 & ITU V.34 & 2400 to 33600 bit/s, after V.8 \\
90 & ITU V.90 & as the digital modem, after V.8 \\
91 & ITU V.91 & two digital modems, after V.8 \\
92 & ITU V.92 & V.90 with quick connect, PCM upstream and
  modem-on-hold \\
56 & K56flex & not implemented: runs V.90 \\
\bottomrule
\end{tabular}
\end{center}

V.90 and V.92 are there as the \emph{digital} modem (the server's side,
as on a G.711 VoIP line), on an 8000 Hz G.711 line: analogue V.90 and
V.92 modems calling connect at up to 56000 bit/s downstream.  90 and 92
offer V.91 too, which two digital modems use instead, and V.34 for the
others; on a line at another sample rate they are V.34.  With automode
V.8 and the V.34 modem hand over to V.32bis (and its V.22bis and V.22)
for a modem without V.34, and the answering modem tries V.23 and Bell
103 for one without V.8.
\end{command}

\begin{command}{AT+MS?}
The settings: \cmd{92,1,300,56000}.
\end{command}

\begin{command}{AT+MS=?}
What may be chosen:\\
\cmd{(21,22,23,122,32,132,34,103,212,56,90,91,92),(0,1),(300-56000),(300-56000)}
\end{command}

\subsection{Identification and country}

\begin{command}{AT+GMI, AT+GMM, AT+GMR}
The manufacturer (\cmd{Smart Link Ltd.}), the model (\cmd{SmartLink
Soft Modem}) and the revision (\cmd{2.9.11}).
\end{command}

\begin{command}{AT+GCI=\p{h}, AT+GCI?, AT+GCI=?}
The country by its T.35 code in hexadecimal (\cmd{+GCI:b5}: USA),
which sets its dial pause (S8) and hook flash (S24) times; \cmd{=?}
lists the codes.  slmodemd's \cmd{-c} option names the country it starts
with (\cmd{--countrylist} lists them).
\end{command}

\subsection{Mode and flow control}

\begin{command}{AT+FCLASS=\p{n}, AT+FCLASS?, AT+FCLASS=?}
\cmd{0} data (default), \cmd{1} fax class 1, \cmd{8} voice
(\cmd{=?} gives \cmd{0,1,8}).
\end{command}

\begin{command}{AT+IFC=2,2}
Hardware flow control both ways: taken, as the pty does its own.  Any
other value is ERROR.
\end{command}

\subsection{V.92 modem-on-hold}

\begin{command}{AT+PMH=\p{n}, AT+PMH?, AT+PMH=?}
\cmd{0} lets the other modem put the call on hold, \cmd{1} refuses
(default 0; \cmd{+PMH?} gives \cmd{+PMH: 0}).
\end{command}

\begin{command}{AT+PMHT=\p{n}, AT+PMHT?, AT+PMHT=?}
How long a hold the modem grants: the code of V.92's Table~33, \cmd{0}
to \cmd{13} (\cmd{0} grants none; default 5).
\end{command}

\begin{command}{AT+PMHR[=\p{s}]}
In online command mode on a V.92 connection: asks the other modem to
go on hold, for \p{s} seconds (1 to 960; 10 if left out).  ERROR
otherwise.
\end{command}

\subsection{Fax class 1}

With \cmd{AT+FCLASS=1}, the modem takes T.31's class 1 commands, the
fax machine's side being the DTE's (the PBX modules and libdmodem have
a T.30 engine of their own as well):

\begin{command}{AT+FTS=\p{n}, AT+FRS=\p{n}}
Sends silence for, or waits for silence of, \p{n} times 10 ms.
\end{command}

\begin{command}{AT+FTM=\p{m}, AT+FRM=\p{m}, AT+FTM=?, AT+FRM=?}
Sends or receives data at modulation \p{m}: CONNECT, then the data with
\cmd{<DLE>} doubled, ended by \cmd{<DLE><ETX>}.
\begin{center}
\begin{tabular}{rl}
\toprule
\p{m} & \\
\midrule
24, 48 & V.27ter, 2400 and 4800 bit/s \\
72, 96 & V.29, 7200 and 9600 bit/s \\
73, 97, 121, 145 & V.17, 7200 to 14400 bit/s, long training \\
74, 98, 122, 146 & V.17, 7200 to 14400 bit/s, short training \\
\bottomrule
\end{tabular}
\end{center}
\end{command}

\begin{command}{AT+FTH=3, AT+FRH=3, AT+FTH=?, AT+FRH=?}
Sends or receives HDLC frames at V.21 channel 2 (300 bit/s), the only
one: \cmd{=?} gives \cmd{3}.
\end{command}

\subsection{Voice}

With \cmd{AT+FCLASS=8} the voice commands of Smart Link's slmodemd are
kept (ERROR in the other modes):

\begin{longtable}{>{\ttfamily}l p{0.62\textwidth}}
\toprule
\normalfont command & \\
\midrule
+VIP & the voice settings as at the start \\
+VNH=\p{n} & 0 or 1 \\
+VGT=\p{n}, +VGR=\p{n} & the gain sending and receiving, 0 to 128 \\
+VIT=\p{n} & the inactivity timer, 0 to 255 \\
+VSD=\p{s},\p{t} & silence detection: sensitivity and period, 0 to 255 \\
+VSM=\p{m},8000 & the compression method, 0 or 1, at 8000 samples/s \\
+VLS=\p{n} & 1 takes the line for voice, 0 leaves it \\
+VTX, +VRX, +VTR & sends, receives, or both \\
+VTS=[\p{f1},\p{f2},\p{d}] & a tone of \p{f1} and \p{f2} Hz for \p{d}
  (0 to 6000) \\
+VTS=\{\p{s},\p{d}\}, +VTS=\p{s} & the DTMF digit \p{s} (for \p{d}, 0 to
  100) \\
+VCID=\p{n} & as \cmd{\#CID} \\
\bottomrule
\end{longtable}
Each takes \cmd{?} (its values) and \cmd{=?} (their ranges).

\section{S-registers}\label{sec:sregs}

Registers 0 to 255 hold 8 bits each.  Those not listed are 0 and free;
``kept'' means set and read back, with no other effect.

\begin{longtable}{rp{0.42\textwidth}rp{0.30\textwidth}}
\toprule
S & & default & \\
\midrule
\endhead
0 & rings to answer after (0: never) & 0 & \\
1 & rings so far & 0 & \\
2 & the escape character & 43 (\cmd{+}) & \\
3 & the line's end & 13 (CR) & \\
4 & the line feed & 10 (LF) & \\
5 & backspace & 8 (BS) & \\
6 & wait for a dial tone, s & 2 & \cmd{X2}, \cmd{X4}; 2 s at least \\
7 & wait for the answer and the training, s & 60 & \\
8 & a comma's pause, s & 2 & the country's \\
9 & carrier detect time, 0.1 s & 6 & kept \\
10 & carrier loss before hanging up, 0.1 s & 7 & \\
11 & DTMF tone and gap, ms & 100 & below 40: 95 \\
12 & escape guard time, 20 ms & 50 & kept: the guard is 0.5 s \\
13 & echo (\cmd{E}) & 1 & \\
14 & quiet (\cmd{Q}) & 0 & \\
15 & verbose (\cmd{V}) & 1 & \\
16 & tone (1) or pulse (0) dialing (\cmd{T}, \cmd{P}) & 1 & \\
24 & hook flash time, 10 ms & 20 & the country's \\
28 & pulse ratio (\cmd{\&P}) & 0 & kept \\
29 & speaker control (\cmd{M}) & 1 & \\
30 & speaker volume (\cmd{L}) & 3 & \\
31 & automode (\cmd{N}, \cmd{+MS}) & 1 & \\
32 & the modulation (\cmd{+MS}) & 92 & \\
39 & line quality (\cmd{\%E}) & 0 & kept \\
56 & results and dial tone (\cmd{X}) & 4 & \\
60 & carrier detect (\cmd{\&C}) & 0 & kept \\
62 & flow control (\cmd{\&H}) & 0 & kept \\
70 & CONNECT's lines (\cmd{\&A}) & 0 & \\
71 & CONNECT's rate (\cmd{\&E}) & 0 & kept \\
92 & the answer's delay, s (0: 10 ms) & 2 & before the answer tone \\
103 & error correction (\cmd{\bs N}) & 1 & \\
104 & data compression (\cmd{\%C}, \cmd{\&K}) & 3 & \\
161 & \cmd{Y} & 0 & kept \\
\bottomrule
\end{longtable}

\section{Online, the escape and hanging up}\label{sec:online}

Online, what the DTE sends goes to the other modem, except the escape:
three S2 characters (\cmd{+++}), each within 0.5 s of the one before,
with 0.5 s of nothing before them and after them.  The modem then says
OK and is in online command mode, the call still up: \cmd{ATO} goes back
online, \cmd{ATH} hangs up, and on a V.92 connection \cmd{AT+PMHR} asks
for modem-on-hold.

The call ends with NO CARRIER when the other side hangs up or its
carrier is gone for S10 tenths of a second.  slmodemd hangs up as well
when the DTE drops DTR (it sets the pty's speed to 0) or, with HUPCL set
on the pty, closes it; it then makes a new pty (\cmd{/dev/ttySL0} follows
it).

\section{slmodemd on a SIP line}\label{sec:sip}

\begin{itemize}
\item \cmd{ATD}\p{number} calls \cmd{sip:}\p{number}\cmd{@}\p{domain} of
  \cmd{SIP\_LOGIN}.  A telephone number (only V.250's dial characters:
  digits, \cmd{* \# + A B C D}, the modifiers \cmd{T P W , @ ! ;} and
  \cmd{- ( ) .} and blanks) goes as its digits and DTMF characters:
  \cmd{ATDT555-1000} calls \cmd{sip:5551000@}\p{domain}.  Anything else
  is a SIP user's name, as it is: \cmd{ATDalice}.
\item The line has no dial tone: \cmd{ATX3} (or \cmd{X1}), else
  NO DIALTONE.
\item With \cmd{-L} slmodemd registers and takes calls: RING every 3
  seconds while one waits, \cmd{ATA} or S0 rings (\cmd{ATS0=1}) answer
  it, and the rings stop if the caller gives up.  A call that comes while
  the line is in use is refused (486, busy); without \cmd{-L} every call
  is (480).
\item The call is G.711 at 8000 Hz, which is what V.90 and V.92 need
  (\cmd{AT+MS=92}, the default).  Over a real PSTN gateway a lower
  modulation (\cmd{AT+MS=132,0,4800,9600}) may do better.
\end{itemize}

\section{Differences from Smart Link's ST56 command set}\label{sec:differences}

\begin{itemize}
\item No stored profiles: \cmd{\&W} and \cmd{\&V} are ERROR, \cmd{\&F}
  and \cmd{Z}\p{n}'s number do nothing; \cmd{ATZ} gives the defaults.
\item No MNP (classes 2 to 5, 10, 10EC): \cmd{\bs A} is taken and
  \cmd{\bs N1} to \cmd{\bs N5} are all V.42; Rockwell's MNP 10 and
  10EC commands (\cmd{-K}, \cmd{-SEC}, \cmd{)M}, \cmd{@M}, \cmd{*H},
  \cmd{-Q}) are ERROR.  A modem with MNP 10EC connects with V.42 (its
  ``EM'' detection patterns are passed over).  See \texttt{docs/link.md},
  MNP.
\item K56flex is not there: \cmd{+MS=56} runs V.90.
\item V.90 and V.92 as the digital modem, V.91, V.44 and V.92's
  modem-on-hold (\cmd{+PMH}, \cmd{+PMHT}, \cmd{+PMHR}) are D-Modem's.
\item \cmd{*B}, \cmd{\bs A} and \cmd{+M}\p{x} give their own result
  (OK, or ERROR for \cmd{+M}\p{x} other than \cmd{+MS}), not the command
  before them.
\item Kept settings that change nothing on a pty or a SIP line:
  \cmd{\&C}, \cmd{\&E}, \cmd{\&H}, \cmd{\%E}, \cmd{\&P}, \cmd{L},
  \cmd{M}, \cmd{Y}; S9 and S12.
\end{itemize}

\end{document}
